\subsection{Substitutions}

Let K be a commutative field, E a unital associative K-algebra, \(\mathbf{x} = (x_i)_{i \in I}\) a family of pairwise permutable elements of E. Let \(\mathbf{B} = \mathbf{K}[(\mathbf{X}_i)_{i \in I}]\) and \(S_x\) the set of all non-zero \(v \in B\) such that \(v(\mathbf{x})\) is invertible in E. Let \(u \in B\) , \(v \in S_x\) and \(f = \frac{u}{v} \in \mathbf{K}((\mathbf{X}_i)_{i \in I})\) . The element \(u(\mathbf{x}) v(\mathbf{x})^{-1} = v(\mathbf{x})^{-1} u(\mathbf{x})\) is defined in E; moreover, if \(u_1\) , \(v_1\) are two polynomials such that \(f = \frac{u_1}{v_1}\) and \(v_1 \in S_x\) , then \(uv_1 = vu_1\) , hence \(u(\mathbf{x}) v_1(\mathbf{x}) = v(\mathbf{x}) u_1(\mathbf{x})\) and so

\[
u (\mathbf {x}) v (\mathbf {x}) ^ {- 1} = u _ {1} (\mathbf {x}) v _ {1} (\mathbf {x}) ^ {- 1}.
\]

Let \(f \in \mathbf{K}((\mathbf{X}_i)_{i \in I})\) . If there exists at least one couple \((u, v)\) such that \(f = \frac{u}{v}\) and \(v \in S_x\) , we shall say that \(x\) is substitutable in \(f\) ; the element \(u(x) v(x)^{-1}\) which only depends on \(f\) and \(x\) is then denoted by \(f(x)\) or \(f((x_i))\) or \(f((x_i)_{i \in I})\) .

PROPOSITION 2. --- Let K be a commutative field, E a unital associative K-algebra and \(\mathbf{x} = (x_i)_{i \in I}\) a family of pairwise permutable elements of E. The set \(\mathbf{S}_{\mathbf{x}}^{-1}\mathbf{B}\) of \(f \in \mathbf{K}((\mathbf{X}_i)_{i \in I})\) such that \(\mathbf{x}\) is substitutable in \(f\) is a K-subalgebra of \(\mathbf{K}((\mathbf{X}_i)_{i \in I})\) . The mapping \(f \mapsto f(\mathbf{x})\) is a unital homomorphism \(\varphi\) of \(\mathbf{S}_{\mathbf{x}}^{-1}\mathbf{B}\) into E. The image \(\varphi(\mathbf{S}_{\mathbf{x}}^{-1}\mathbf{B})\) is the set of all \(yz^{-1}\) , where \(y\) runs over the unital subalgebra \(\mathbf{K}[\mathbf{x}]_E\) of E generated by the family \(\mathbf{x}\) and where \(z\) runs over the set of all invertible elements of \(\mathbf{K}[\mathbf{x}]_E\) .

Let \(f_{1} = \frac{u_{1}}{v_{1}}\) , \(f_{2} = \frac{u_{2}}{v_{2}}\) be two elements of \(K((X_{i})_{i\in I})\) such that \(v_{1}, v_{2} \in S_{x}\) . We have \(f_{1} + f_{2} = \frac{u_{1}v_{2} + u_{2}v_{1}}{v_{1}v_{2}}\) , \(f_{1}f_{2} = \frac{u_{1}u_{2}}{v_{1}v_{2}}\) and \(v_{1}, v_{2} \in S_{x}\) . Hence \(S_{x}^{-1}B\) is a K-subalgebra of \(K((X_{i})_{i\in I})\) . The rest of the proposition is clear.

COROLLARY. --- Let L be a commutative field, K a subfield of L, \(\mathbf{x} = (x_i)_{i \in I}\) a family of elements of L, M the set consisting of the \(x_i\) , U the set of all \(f \in \mathrm{K}((\mathrm{X}_i)_{i \in I})\) such that x is substitutable in f and \(\varphi\) the homomorphism \(f \mapsto f(\mathbf{x})\) of U into L. Then \(\varphi(\mathrm{U})\) is the subfield of L generated by \(K \cup M\) .

Let \(L'\) be the subfield of L generated by \(K \cup M\) . We have

\[
\mathbf {K} \cup \mathbf {M} \subset \varphi (\mathbf {U}) \subset \mathbf {L} ^ {\prime}
\]

and \(\varphi(\mathbf{U})\) is a subring of \(\mathbf{L}\) . Now Prop. 2 implies that \(\varphi(\mathbf{U})\) is a subfield of \(\mathbf{L}\) , whence \(\varphi(\mathbf{U}) = \mathbf{L}'\) .

Let \(f \in \mathbf{K}((\mathbf{X}_i)_{i \in I})\) and let \((g_i)_{i \in I}\) be a family of elements of \(\mathbf{K}((\mathbf{Y}_l)_{l \in L})\) . If \((g_{i})\) is substitutable in f, then \(f((g_{i}))\) is an element of \(\mathbf{K}((\mathbf{Y}_{l})_{l\in L})\) . In particular, \((\mathbf{X}_{i})_{i\in I}\) is substitutable in f and \(f=f((\mathbf{X}_{i})_{i\in I})\) .

PROPOSITION 3. --- Let E be an algebra over K which is associative, commutative, unital and non-zero. Let \(f \in \mathbf{K}((\mathbf{X}_i)_{i \in \mathbb{I}})\) and for each \(i \in \mathbb{I}\) , let \(g_i \in \mathbf{K}((\mathbf{Y}_l)_{l \in \mathbb{I}})\) . Given a family \(\mathbf{y} = (y_l)_{l \in \mathbb{L}}\) of elements of E, suppose that \(\mathbf{y}\) is substitutable into each \(g_i\) and \((g_i(\mathbf{y}))_{i \in \mathbb{I}}\) is substitutable in \(f\) . Then:

\begin{enumerate}
\def\labelenumi{(\roman{enumi})}
\item
  \((g_{i})_{i\in I}\) is substitutable in \(f\) ;
\item
  if we denote by \(h\) the element \(f((g_i))\) of \(\mathbf{K}((\mathbf{Y}_l)_{l \in L})\) , then \(\mathbf{y}\) is substitutable in \(h\) and \(h(\mathbf{y}) = f((g_i(\mathbf{y})))\) .
\end{enumerate}

We may take I to be finite. By hypothesis, for each \(i \in I\) , \(g_i\) can be put in the form \(p_i/q_i\) where \(p_i\) , \(q_i \in K[(Y_l)_{l \in L}]\) and \(q_i(\mathbf{y})\) is invertible in E. Likewise \(f\) can be written as \(u/v\) , where \(u\) , \(v \in K[(X_i)_{i \in I}]\) and \(v((g_i(\mathbf{y})))\) is invertible. Let \(m = \sup(\deg u, \deg v)\) , and let \(w = \prod_{i \in I} q_i \in K[(Y_l)_{l \in L}]\) , \(u_1 = u((g_i))w^m\) , \(v_1 = v((g_i))w^m\) . The polynomial \(u\) is a K-linear combination of monomials \(\prod_{i \in I} X_i^{\nu_i}\) such that \(\sum_{i \in I} \nu_i \leqslant m\) . We have \(w^m \prod_{i \in I} g_i^{\nu_i} = w^m\left(\prod_{i \in I} p_i^{\nu_i}\right)\left(\prod_{i \in I} q_i^{\nu_i}\right)^{-1} \in K[(Y_l)_{l \in L}]\) by the choice of \(m\) . Hence \(u_1 \in K[(Y_l)_{l \in L}]\) and similarly \(v_1 \in K[(Y_l)_{l \in L}]\) . Moreover, \(v_1(\mathbf{y}) = (w(\mathbf{y}))^m v((g_i(\mathbf{y})))\) is invertible. Hence \(v_1 \neq 0\) , because \(E \neq 0\) , and so \(v((g_i)) \neq 0\) . The family \((g_i)\) is thus substitutable in \(f\) . Besides we have \(f((g_i)) = u_1/v_1\) , hence \(\mathbf{y}\) is substitutable in \(h = f((g_i))\) , and \(h(\mathbf{y}) = u_1(\mathbf{y})/v_1(\mathbf{y}) = u((g_i(\mathbf{y})))/v((g_i(\mathbf{y}))) = f((g_i(\mathbf{y})))\) .

Let K be a commutative field, E a commutative associative and unital K-algebra, and let \(f \in \mathbf{K}((\mathbf{X}_i)_{i \in I})\) . Let \(T_f\) be the set of all \(\mathbf{x} = (x_i)_{i \in I} \in E^{I}\) which are substitutable in \(f\) . The mapping \(\mathbf{x} \mapsto f(\mathbf{x})\) of \(T_f\) into E is called the rational function associated with \(f\) (and E); we sometimes denote it by \(\tilde{f}\) . If \(g \in \mathbf{K}((\mathbf{X}_i)_{i \in I})\) we have \(T_f \cap T_g \subset T_{f + g}\) , \(T_f \cap T_g \subset T_{fg}\) , hence the rational function associated with \(f + g\) (resp. \(fg\) ) is defined on \(T_f \cap T_g\) and has the same value on this set as \(\tilde{f} + \tilde{g}\) (resp. \(\tilde{f}\tilde{g}\) ). Let \(T_f'\) be the set of \(\mathbf{x} \in T_f\) such that \(f(\mathbf{x})\) is invertible; if \(\mathbf{x} \in T_f'\) , \(\mathbf{x}\) is substitutable in \(1/f\) and the rational function associated with \(1/f\) takes at \(\mathbf{x}\) the value \(f(\mathbf{x})^{-1}\) .

If K is an infinite commutative field, \(f \in \mathbf{K}((\mathbf{X}_i)_{i \in I})\) , \(g \in \mathbf{K}((\mathbf{X}_i)_{i \in I})\) and \(\tilde{f}\) , \(\tilde{g}\) are the rational functions associated with f, g (and K), and if \(\tilde{f}(\mathbf{x}) = \tilde{g}(\mathbf{x})\) for all \(x \in T_f \cap T_g\) then f = g. For if f = u/v and \(g = u_1 / v_1\) , where u, v, \(u_1\) , \(v_1\) are polynomials, we have \(u(\mathbf{x})v_1(\mathbf{x}) = u_1(\mathbf{x})v(\mathbf{x})\) for all x such that \(v(x)v_1(x) \neq 0\) , hence \(uv_1 = u_1v\) (IV, p.~18, Th. 2). Therefore the mapping \(f \mapsto \tilde{f}\) is injective and we shall often identify f and \(\tilde{f}\) .

* Using the factoriality of \(\mathbf{K}[(\mathbf{X}_i)_{i\in I}]\) (Comm. Alg., VII, § 3, No.~2 p.~502 and Cor. of Th. 2 p.~506), one easily shows the following: for every \(f\in \mathbf{K}((\mathbf{X}_i)_{i\in I})\) there exist \(u,v\in \mathbf{K}[(\mathbf{X}_i)_{i\in I}]\) such that:

\begin{enumerate}
\def\labelenumi{\arabic{enumi})}
\item
  \(f = u / v\)
\item
  for \(\mathbf{x} \in \mathbf{K}^{I}\) to be substitutable in \(f\) it is necessary and sufficient that \(v(\mathbf{x}) \neq 0\) .
\end{enumerate}

